% \ifcat and \if -- catcodes and character codes, both *expanding*.
%
% \ifcat was false for everything, \ifcat\relax\relax included: it compared
% only (Char, Char) and never expanded its operands. The rule was written in
% the comment directly above the code and simply not implemented beneath it.
%
% tex.web 506: a control sequence is compared as category code 16 and
% character code 256, so any two control sequences are equal under both
% conditionals.
%
% expl3 ends every floating-point expression with this test:
%
%   \if_catcode:w \scan_stop: \exp_not:N #2
%
% which asks "is #2 a control sequence" -- how \__fp_parse_infix:NN recognises
% its end marker. Answering "no" made l3fp read the marker as another operand,
% i.e. implicit multiplication, so \c_e_fp came out as `e * <end marker>`: an
% invalid operation, reported in a loop that never terminated.
%
% \noexpand marks the token that follows, and the marker is never itself an
% operand -- \exp_not:N sits in front of every token these conditionals see.
\input plain
\def\M#1{\message{<#1>}}
\def\bar{x}
\let\mark=\relax
\chardef\ch=65
\M{a=\ifcat\relax\noexpand\undefinedcs T\else F\fi}
\M{b=\ifcat\relax\noexpand a T\else F\fi}
\M{c=\ifcat\relax\relax T\else F\fi}
\M{d=\ifcat\relax\noexpand\bar T\else F\fi}
\M{e=\ifcat\noexpand\bar\noexpand\relax T\else F\fi}
\M{f=\ifcat\relax\noexpand\mark T\else F\fi}
\M{g=\ifcat\noexpand\mark\noexpand\bar T\else F\fi}
\M{h=\ifcat\relax\noexpand+ T\else F\fi}
\M{i=\ifcat\noexpand a\noexpand b T\else F\fi}
\M{j=\ifcat\relax\noexpand\ch T\else F\fi}
\M{k=\ifcat ab T\else F\fi}
\M{l=\ifcat a+ T\else F\fi}
\M{m=\ifcat +- T\else F\fi}
\M{n=\ifcat \bar a T\else F\fi}
\def\two{ab}
\M{o=\ifcat \two b T\else F\fi}
\M{p=\if ab T\else F\fi}
\M{q=\if aa T\else F\fi}
\M{r=\if\relax\relax T\else F\fi}
\M{s=\if\noexpand\bar\noexpand\mark T\else F\fi}
\M{t=\ifcat\noexpand\ch\noexpand\bar T\else F\fi}
\bye
